\section{Julia effect estimators and the meaning of a reusable result}
\label{exec21:effects-scope}

A servicing manager asks whether an additional contact is worth its cost for
each risk group. A numerical answer requires more than a treatment-effect
method name. The population average, the average among previously contacted
customers, a conditional contrast and a coefficient from residual regression
can differ even when every conditional model is known. Missing outcomes can
change the averaging population again. A reusable interface must carry this
meaning from the business question through the numerical calculation.

The selected Julia sources offer several distinct building blocks.
CausalELM~0.8.0 supplies g-computation and effect learners; CausalForest~0.1.5
implements tree-based contrasts; TMLE~0.20.5 supplies targeted and one-step
estimation; and CausalTables~1.3.4 represents causal data and interventions in
supplied structural models. Their selected source revisions precede the
7~October~2026 cutoff. Their MIT notices permit substantial reuse subject to
retaining the specified notices, but do not establish dependency clearance or
statistical validity. The following comparison concerns inspected source
paths, not native execution or a tested combination of these packages.
\cite{exec21_effects_causalelm,exec21_effects_causalforest,exec21_effects_tmle,exec21_effects_causaltables}

\subsection{An average effect is a population-specific quantity}
\label{exec21:effects-target}

Let $W$ denote pre-contact risk information, $A\in\{0,1\}$ the contact action,
and $Y(a)$ the monetary outcome under action $a$. Under consistency,
conditional exchangeability and treatment positivity, write
$e(w)=\Pr(A=1\mid W=w)$ and
$\tau(w)=E\{Y(1)-Y(0)\mid W=w\}$. The population average treatment effect is
$E\{\tau(W)\}$; the treated-population effect is
$E\{\tau(W)\mid A=1\}$. A prediction vector containing estimated
$\tau(W_i)$ is useful, but does not by itself provide uncertainty for a
selected subgroup or a policy chosen using that vector.

In the inspected CausalELM g-formula path, ATE, ITT and CATE labels use
counterfactual predictions at all supplied covariate rows; the other branch
uses treated rows. The g-computation wrapper averages the contrast, while the
S-learner retains a contrast vector. The T-learner fits each treatment arm and
subtracts its predictions on the supplied covariates. Thus an ITT label only
has an assignment-effect interpretation when the supplied treatment actually
represents assignment; it cannot convert receipt of contact into randomized
offer assignment. Categorical coding, the action clock and the averaging
population must be explicit inputs. \cite{exec21_effects_causalelm}

The same source's double-machine-learning path forms outcome and treatment
residuals and averages foldwise least-squares coefficients. Its fold
constructor slices contiguous rows. Neither the name nor contiguous slicing
supplies a dependence argument for repeated customer records. Even with
oracle nuisance functions and independent observations, a partially linear
coefficient need not be the population average under heterogeneous effects.
To see the distinction, take
$Y=m_0(W)+A\tau(W)+\varepsilon$ with
$E(\varepsilon\mid A,W)=0$. If $m(W)=E(Y\mid W)$, then
$Y-m(W)=(A-e(W))\tau(W)+\varepsilon$. Consequently the population residual
coefficient is
\begin{equation}
\label{exec21:effects-residual}
 \theta_{\rm res}
 =\frac{E[(A-e(W))(Y-m(W))]}{E[(A-e(W))^2]}
 =\frac{E[e(W)(1-e(W))\tau(W)]}{E[e(W)(1-e(W))]}.
\end{equation}
This is an overlap-weighted average in this model. A constant effect or
additional conditions can make it agree with the ATE; cross-fitting alone
does not do so. This derivation concerns the oracle population target and is
not a claim that a fitted implementation attains it in finite samples.
\cite{exec21_effects_causalelm,exec21_effects_reference}

Return to the stipulated servicing population used for the R forest in
Section~\ref{exec21:grfr-target}: the two risk groups have equal population
mass, effects $(-4,12)$ dollars and treatment propensities $(1/2,1/10)$.
Using the same law makes the target comparison across languages explicit;
it supplies no independent empirical corroboration. Exact enumeration gives
\begin{equation}
\label{exec21:effects-targetexample}
 \tau_{\rm ATE}=4,\qquad
 \tau_{\rm ATT}=-\frac43,\qquad
 \theta_{\rm res}=\frac4{17}.
\end{equation}
At an assumed cost of USD~2 per additional contact, the population margin is
USD~2. Subtracting that cost from the residual coefficient instead gives
$-30/17$ dollars. These quantities answer different questions. A proposed
product should refuse a request to relabel the coefficient as the
population-wide incremental value. A separately evaluated targeted policy
could be more useful, but selecting the profitable subgroup and estimating
its value require their own information and uncertainty contract.
\cite{exec21_effects_reference}

\subsection{Source-level uncertainty checks before adoption}
\label{exec21:effects-uncertainty}

The inspected CausalELM generic null routine samples treatment independently
and uniformly from its distinct observed values, or uniformly over its
observed range for a continuous treatment. This does not preserve the
observed treatment count. It also does not reproduce a nonuniform propensity
law or a cluster-randomized design without additional construction. For four
customers assigned exactly one contact, independent fair binary draws retain
that design with probability $4/16=1/4$ and create an empty treatment or
control arm with probability $2/16=1/8$. This exact calculation identifies an
assignment-law mismatch, rather than estimating a package error rate.
\cite{exec21_effects_causalelm,exec21_effects_reference}

There is also a direct dimensional problem in the selected source. For null
statistics $z_1,\ldots,z_B$ and reported statistic $t$, its scalar
\texttt{stderr} expression is
\begin{equation}
\label{exec21:effects-scale}
 s_{\rm code}=\frac{\sum_{b=1}^{B}(t-z_b)^2}{(B-1)\sqrt B}.
\end{equation}
If $t$ and $z_b$ are monetary effects, this expression has squared monetary
units. Multiplication of all outcomes by a positive constant $c$ multiplies
it by $c^2$, whereas a standard error must scale by $c$. With
$(z_1,z_2,z_3,z_4)=(-1,-1,1,1)$ and $t=2$, the expression is $10/3$.
Changing dollars to cents multiplies it by 10,000 rather than 100. This is a
source-level defect in the expression at the selected revision. It does not
require a Julia run to establish the algebra, and it does not establish the
finite-sample behavior of every other package function.
\cite{exec21_effects_causalelm,exec21_effects_reference}

Taking a square root in one location would not settle the larger inference
question. The sampled null law, the centering, the effect being tested and
the relationship between repeated randomization draws and sampling
uncertainty would still need justification. The inspected interval routine
searches a finite grid with tail comparisons against the same null draws;
that code alone is insufficient to claim a valid general confidence-interval
inversion. The proposed reuse decision is therefore to reject unqualified
reuse of this uncertainty path, while retaining the point-estimation family
as a candidate for a separately justified implementation.
\cite{exec21_effects_causalelm}

CausalForest illustrates a different boundary. Its interface exposes
bootstrap or subsampling choices and separate split and leaf samples under
the honest-tree option. However, its outcome-centering forest is fitted
before these per-tree splits. The inspected prediction paths form
within-leaf treated and control weights and return centered-outcome
contrasts. They return point predictions, rather than a calibrated covariance
object. Honesty of the tree split is consequently not sufficient evidence
that all nuisance fitting is independent of evaluation, nor evidence of
equivalence to the separately reviewed \texttt{grf} implementation. A native
comparison should examine the actual nuisance training rows, treatment
coding, arm support in leaves, random-number state and failure behavior
before considering any shared inference contract.
\cite{exec21_effects_causalforest}

\subsection{Targeted estimation, missingness and influence scales}
\label{exec21:effects-tmle}

TMLE exposes nuisance-model choices, outer resampling and propensity
truncation. The selected noncollaborative default uses no outer resampling;
a requested cross-validation design must be supplied and preserved. The
implementation's \texttt{weighted} constructor default is true although
adjacent prose says false, another reason to record actual resolved
configuration. Truncation prevents numerical division by a very small
estimated probability but cannot create absent treatment support. For an
identified ATE in an independent sampling model, the familiar influence
expression clarifies what an adapter must carry:
\begin{equation}
\label{exec21:effects-influence}
\begin{aligned}
 D_i={}&\frac{A_i}{e(W_i)}\{Y_i-m_1(W_i)\}
       -\frac{1-A_i}{1-e(W_i)}\{Y_i-m_0(W_i)\}\\
       &+m_1(W_i)-m_0(W_i)-\psi .
\end{aligned}
\end{equation}
The relevant source combines outcome-residual and counterfactual-mean
components. It stores the standard deviation of the estimated influence
values together with their number, and passes both to the testing
interface. This scale is not already the standard error of the effect.
Conditional asymptotic inference additionally needs suitable nuisance
remainder control, positivity, moments and a sampling-law argument; the
presence of an influence vector does not prove those conditions.
\cite{exec21_effects_tmle}

For example, influence values $(-3,-1,1,3)$ have sample variance $20/3$.
Under the independent mean calculation, the variance of the average is
$5/3$, since $n=4$. Exporting the influence standard deviation as an effect
standard error would miss the factor $\sqrt n$. Conversely, dividing a
quantity that is already a standard error by $\sqrt n$ again would understate
uncertainty. For joint outputs, the source stacks influence vectors and
computes their covariance. A bridge must preserve the identity and order
of contributing customers; equal vector length is insufficient when
different outcomes select different complete cases. Cluster or calendar
dependence requires an appropriate covariance construction, not a new field
name on the independent calculation. \cite{exec21_effects_tmle,exec21_effects_reference}

The missingness documentation is explicit: rows missing variables needed
for estimation are dropped. Without cross-validation, propensity fitting
may use additional records with observed treatment and confounders even
when outcomes are missing. The source's final fluctuation dataset nevertheless
uses complete relevant records. This is useful data reuse, but is not a
general missing-outcome correction. \cite{exec21_effects_tmle}

Let $R$ indicate observation of the outcome. Consider independent customers
with equally frequent risk groups, randomized contact probability $1/2$,
effects $(-4,12)$ and outcome-response probabilities $(9/10,1/10)$ by risk
group. Suppose response is independent of potential outcomes and contact
conditional on risk, so missing at random holds in this example and response
positivity holds. The complete cases have group probabilities $(9/10,1/10)$,
giving
\begin{equation}
\label{exec21:effects-response}
 E\{\tau(W)\}=4,\qquad
 E\{\tau(W)\mid R=1\}=-\frac{12}{5}.
\end{equation}
Correct conditional response weighting recovers 4 in the exact known law.
At cost USD~2, the population margin is USD~2 and the complete-case margin
is USD~$-22/5$. The example does not claim these are observed package
outputs. It shows why improving a treatment-propensity fit cannot, by
itself, repair a changed marginal population. A missing-data-aware
identification and estimation design is required for the original request.
\cite{exec21_effects_reference}

Case-control sampling is also a distinct design. The inspected TMLE
interface accepts a population prevalence as a scalar and constructs
prevalence weights. Uncertainty in an externally estimated prevalence is
not an input to those inspected interfaces. A banking adapter must therefore
state whether prevalence is assumed fixed or estimated and how its
uncertainty enters a proposed joint analysis. It should not silently reuse
case-control weights for nonresponse, population transport or clustered
sampling, whose information and identifying assumptions differ.
\cite{exec21_effects_tmle}

\subsection{Supplied causal models are valuable public references}
\label{exec21:effects-scm}

CausalTables supplies a useful route to inexpensive comparisons before bank
experiments. Its intervention functions modify treatment, recompute
summaries and draw counterfactual responses from a supplied structural
causal model. Its \texttt{cfmean} and \texttt{cfdiff} functions approximate
intervention means by Monte Carlo. Despite efficiency-bound wording in
their documentation, the inspected return tuples contain only the mean.
A reference-generation service must separately retain the generating
equations, intervention, random seed, Monte Carlo error and dependence
structure. A simulated network summary does not establish a variance formula
for observations from a bank network. \cite{exec21_effects_causaltables}

The distinction between supplied truth and observational identification is
not semantic. Let $U$ be a fair binary variable and consider two structural
models. In both, $A=U$. In the first, $Y=A$; in the second, $Y=U$. Both
produce the same observed law, with probability $1/2$ on each of
$(A,Y)=(0,0)$ and $(1,1)$. Yet the intervention contrast is 1 in the first
model and 0 in the second. A simulator can reproduce either contrast once
its model is supplied. The observed distribution alone does not choose
between them. At an illustrative cost of $1/2$, the corresponding margins
are $1/2$ and $-1/2$. \cite{exec21_effects_reference}

TMLE's CausalTables extension converts treatment columns to categorical
arrays and identifies an estimand using the causes recorded in the supplied
table. This is an explicit and promising connection, provided that level
order, graph semantics and the provenance of the causal declarations survive
conversion. It does not verify that the declared graph is true. The
matching CausalInference revision supplies separately reviewed graphical
interfaces; that evidence should be reused for its established scope rather
than treated as qualification of heterogeneous-effect estimation.
\cite{exec21_effects_tmle,exec19_otherid_julia}

\subsection{A bounded adoption path and the wider language search}
\label{exec21:effects-integration}

The proposed architecture is an isolated Julia worker with a pinned
environment and an explicit estimand/result schema. TMLE is a candidate for
targeted mean estimation and CausalTables for supplied-model reference
generation. A first public comparison should use exact finite distributions
to verify scalar targets, treatment-level order, outcome-unit conversion,
complete-case selection and graph-to-estimand conversion. Subsequent native
trials should compare influence-vector scaling, fold membership, fitted
nuisance behavior and refusal on absent support. CausalELM's selected
uncertainty path needs a justified repair before inclusion; CausalForest's
point predictions need their own inference qualification. Porting either
path before understanding these contracts would copy its unresolved
scientific meaning into a new language.

The library, CLI and MCP representations should carry the action/assignment
distinction, target population, contributing unit IDs, missingness rule,
fitted-model configuration, uncertainty kind and all unavailable results.
An agent can propose mappings and inspect contradictory declarations, but
must leave unsupported identification or uncertainty requests unresolved.
This gives the business user a meaningful answer when a requested population
effect cannot be supplied, and permits genuine comparison of methods that
do answer the same question.

Rust coverage remains incomplete. Separate official repository and crate
queries identify leads including \texttt{perpetual}, \texttt{causalkit},
\texttt{oxicuda-causal} and \texttt{propensity-score}, alongside graph and
event-ordering projects. These are metadata leads; current descriptions are
not proof of functionality at the cutoff. The queries inspect only their
first pages, and an exact treatment-effect phrase query returning zero does
not establish absence. Selected pre-cutoff source revisions, licenses,
estimators and inference paths must be inspected before any capability or
reuse conclusion. Wider language coverage and package-paper originals also
remain to be established. \cite{exec21_effects_discovery}

These public comparisons can determine which Julia integrations merit a
native prototype and which assumptions a later banking experiment must test.
